% Part: sets-functions-relations
% Chapter: functions
% Section: kinds

\documentclass[../../../include/open-logic-section]{subfiles}

\begin{document}

\olfileid{sfr}{fun}{kin}
\olsection{د تابعو ډولونه}

\begin{explain}
د تابعو د هغو ځينو ډولونو دپاره چې ډېر ورسره مخ کېږو، د يو ډول
ډلبندۍ پېژندل به ګټور وي۔

په پيل کښې داسې تابعې په پام کښې نيول غواړو چې د هدف سټ هر
غړے د تابعې يو قيمت وي۔ دې تابعو ته !!{surjective} ويل کېږي،
او د \olref{fig:surjective} په شان ئې ښودلے شو۔

\begin{figure}
  \olasset{assets/diagrams/surjective.tikz}
  \caption{!!^a{surjective} تابع د هدف سټ هر !!{element} د قيمت په توګه لري۔}
  \ollabel{fig:surjective}
\end{figure}
\end{explain}

\begin{defn}[!!^{surjective} تابع]
تابع $f \colon A \rightarrow B$ هله او يوازې هله \emph{!!{surjective}}
ده چې $B$ د~$f$ د قيمتونو سټ هم وي؛ يعنې د هر $y \in B$
دپاره لږ تر لږه يو $x \in A$ وي چې~$f(x) = y$ وي، يا په نښو:
\[
  (\forall y \in B)(\exists x \in A)f(x) = y.
\]
داسې تابع له $A$ څخه $B$ ته !!a{surjection} بولو۔
\end{defn}

\begin{explain}
که ښودل غواړئ چې $f$ يوه !!a{surjection} ده، نو بايد وښيئ چې د~$f$
د هدف سټ هر څيز د کوم ننوت $x$ دپاره د $f(x)$ قيمت دے۔

پام وکړئ چې هره تابع يوه !!a{surjection} \emph{رامنځته کوي}۔ ځکه،
د ورکړل شوې تابعې $f \colon A \to B$ دپاره، فرض کړئ
$f' \colon A \to \ran{f}$ په $f'(x) = f(x)$ تعريف شوې ده۔ ځکه
چې $\ran{f}$ د $\Setabs{f(x) \in B}{x \in A}$ په توګه \emph{تعريف}
شوے دے، دا تضمين شوے دے چې تابع $f'$ به !!a{surjection} وي۔
\end{explain}

\begin{explain}
اوس، هره تابع هر ممکن ننوت له يوه يوازيني وت سره تړي۔ خو داسې
تابعې هم شته چې بېل ننوتونه هېڅکله له يو شان وت سره نۀ تړي۔ دې
تابعو ته !!{injective} ويل کېږي، او د \olref{fig:injective} په شان
ئې ښودلے شو۔
\begin{figure}
  \olasset{assets/diagrams/injective.tikz}
  \caption{!!^a{injective} تابع هېڅکله دوه بېل ارګومېنټونه له يوه قيمت سره نۀ تړي۔}
  \ollabel{fig:injective}
\end{figure}
\end{explain}

\begin{defn}[!!^{injective} تابع]
تابع $f \colon A \rightarrow B$ هله او يوازې هله \emph{!!{injective}}
ده چې د هر $y \in B$ دپاره تر يوه زيات داسې $x \in A$ نۀ وي
چې~$f(x) = y$ وي۔ داسې تابع له $A$ څخه~$B$ ته !!a{injection} بولو۔
\end{defn}

\begin{explain}
که ښودل غواړئ چې $f$ يوه !!a{injection} ده، بايد وښيئ چې د~$f$
د تعريف د ساحې د هر دوو \psOblique{!!{element}s} $x$ او $y$ دپاره،
که $f(x)=f(y)$ وي، نو $x=y$ وي۔
\end{explain}

\begin{ex}
ثابته تابع $f\colon \Nat \to \Nat$ چې په $f(x) = 1$ ورکړل شوې،
نۀ !!{injective} ده او نۀ !!{surjective}۔

د عينيت تابع $f\colon \Nat \to \Nat$ چې په $f(x) = x$ ورکړل شوې،
هم !!{injective} ده او هم !!{surjective}۔

د ځايناستي تابع $f \colon \Nat \to \Nat$ چې په $f(x) = x+1$ ورکړل
شوې، !!{injective} ده خو !!{surjective} نۀ ده۔
  
تابع $f \colon \Nat \to \Nat$ چې داسې تعريف شوې ده:
\[
  f(x) =
  \begin{cases}
    \frac{x}{2} & \text{که $x$ جفت وي} \\
    \frac{x+1}{2} & \text{که $x$ طاق وي۔}
  \end{cases}
\]
!!{surjective} ده، خو !!{injective} نۀ ده۔
\end{ex}

\begin{explain}
ډېر وخت داسې تابعې په پام کښې نيول غواړو چې هم !!{injective} وي
او هم !!{surjective}۔ دا تابعې !!{bijective} بولو۔ دوې په
\olref{fig:bijective} کښې ښودل شوې تابعې ته ورته دي۔ !!^{bijection}s
کله کله \emph{يو پر يو مطابقتونه} هم بلل کېږي، ځکه چې د هدف سټ
غړي د تعريف د ساحې له غړو سره په يوازيني ډول جوړوي۔
\begin{figure}
  \olasset{assets/diagrams/bijective.tikz}
  \caption{!!^a{bijective} تابع د هدف سټ غړي د تعريف د ساحې له غړو سره په يوازيني ډول جوړوي۔}
  \ollabel{fig:bijective}
\end{figure}
\end{explain}

\begin{defn}[!!^{bijection}]
تابع $f \colon A \to B$ هله او يوازې هله \emph{!!{bijective}} ده چې
هم !!{surjective} وي او هم !!{injective}۔ داسې تابع له $A$ څخه~$B$
ته (يا د $A$ او~$B$ ترمنځ) !!a{bijection} بولو۔
\end{defn}

\end{document}
